% Shared exercise chunk: l2-probability-rules
\begin{exercisechunk}
\item \textbf{Probability rules and their notation.}
\label[exercise]{ex:probability-rules}
Use the definitions in \cref{app:probability}.
\begin{enumerate}[(a)]
\item Derive the following from the probability axioms:
\begin{enumerate}[label=(\roman*),leftmargin=*,beginpenalty=10000]
\item $\P(\emptyset)=0$.
\item $\P(E^c)=1-\P(E)$ for every event $E$.
\item Monotonicity of probability.
\end{enumerate}
\item Events and indicators:
\begin{enumerate}[label=(\roman*),leftmargin=*,beginpenalty=10000]
\item Write $\P(X=x,Y\in D)$ and $\I(X\ne Y)$ with their events expressed
as subsets of $\Omega$.
\item Show that $\I(E)$ is a measurable real-valued
function whenever $E$ is an event.
\end{enumerate}
\item Expectations:
\begin{enumerate}[label=(\roman*),leftmargin=*,beginpenalty=10000]
\item For finite-valued random variables with an arbitrary joint distribution,
prove linearity of expectation directly from the finite-sum definition.

\item Also prove $\mathbb E[\I(E)]=\P(E)$.
\end{enumerate}
\item For a finite partition, prove each rule below, accounting for
zero-probability parts of the partition:
\begin{enumerate}[label=(\roman*),leftmargin=*,beginpenalty=10000]
\item The multiplication rule.
\item Total probability.
\item Total expectation.
\end{enumerate}
\item Conditioning and independence:
\begin{enumerate}[label=(\roman*),leftmargin=*,beginpenalty=10000]
\item Prove the rule in \cref{eq:conditional-block-independence} by writing
conditional probabilities as ratios.
\item Give an example where two independent
random variables become dependent after conditioning on an event involving
both of them.
\end{enumerate}
\end{enumerate}
\end{exercisechunk}
